% !TEX root = ../main.tex

\begin{titlepage}
\begin{abstract}
%
Le API sono oggi il principale punto di accesso ai servizi cloud, e l'OWASP API
Security Top 10 (2023) ne elenca i rischi più critici. Non tutte le dieci
categorie, però, lasciano una traccia negli artefatti di progetto: alcune si
riconoscono nel codice, nel contratto OpenAPI o nella configurazione, altre
dipendono dal comportamento del sistema in esecuzione o dalla logica di business.
Gli strumenti di analisi statica non distinguono questi casi, e una loro
segnalazione, o la sua assenza, non dice quanto la categoria sia davvero
coperta.

La tesi propone un criterio che classifica le dieci categorie in tecniche,
rilevabili interamente con l'analisi statica, ibride, rilevabili solo in parte,
e di logica di business, non rilevabili sugli artefatti. Per rendere il criterio
operativo è stato realizzato uno strumento che combina scanner esistenti
(Checkov, Semgrep e Spectral) con un modulo di rilevamento per ciascuna categoria
tecnica o ibrida, e che affianca all'analisi statica una verifica dinamica dove
questa non basta. Lo strumento è stato validato su coppie di applicazioni
gemelle, una vulnerabile e una corretta, sul benchmark crAPI e su un insieme di
prova di tre applicazioni di terzi: VAmPI, DVAPI e sadcloud.

Quattro categorie risultano tecniche, quattro ibride e una di logica di
business, mentre API9 resta fuori dalla classificazione. Sulle coppie gemelle
ogni modulo distingue la versione vulnerabile da quella corretta senza falsi
positivi, e su VAmPI, che ha lo stesso codice nelle due modalità, solo la
verifica dinamica separa quella vulnerabile da quella sicura, come la
classificazione prevede. Su codice di terzi lo strumento trova il 61\% delle
misconfiguration documentate dell'infrastruttura di sadcloud, ma riconosce
poche vulnerabilità delle API: la BOLA di VAmPI e la SSRF di crAPI. Il divario
dipende soprattutto dalle forme di codice che le regole conoscono, e indica la
direzione in cui estendere lo strumento.
%
\\[1cm]
\end{abstract}
\end{titlepage}
